\lyseis{Β1.2}{Γωνία, γραμμή, σχήματα, ίσα σχήματα}
\sol{1} Κορυφή το $\mathrm{Ο}$, πλευρές οι ημιευθείες $\mathrm{Οx}$ και $\mathrm{Οy}$, εσωτερικό τα σημεία ανάμεσά τους.
\sol{2} Δύο γωνίες: μία κυρτή και μία μη κυρτή.
\sol{3} α) ευθεία (ανοιχτή) \quad β) ανοιχτή τεθλασμένη \quad γ) κλειστή καμπύλη \quad δ) ανοιχτή καμπύλη \quad ε) κλειστή τεθλασμένη. Τεθλασμένες είναι οι β και ε.
\sol{4} Τρίγωνο: $3,\,3,\,0$. Τετράπλευρο: $4,\,4,\,2$. Πεντάγωνο: $5,\,5,\,5$. Εξάγωνο: $6,\,6,\,9$. Οκτάγωνο: $8,\,8,\,20$.
\sol{5} $\mathrm{ΑΓ}$, $\mathrm{ΑΔ}$, $\mathrm{ΒΔ}$, $\mathrm{ΒΕ}$, $\Gamma\mathrm{Ε}$.
\sol{6} Ίσα είναι τα ζεύγη α και γ.
\sol{7} α) Σωστό. β) Σωστό. γ) Λάθος. δ) Λάθος. ε) Σωστό.
\sol{8} $5$ σχοινιά, όσες και οι διαγώνιοι.
\sol{9} Ελεύθερη απάντηση: ίσα είναι τα πλακάκια που εφαρμόζουν ακριβώς.
\sol{10} Ανοιχτή τεθλασμένη. Αν ενώσουμε το $\mathrm{Δ}$ με το $\mathrm{Α}$, παίρνουμε τετράπλευρο.
\sol{11} Τέσσερις γωνίες.
\sol{12} Ελεύθερη απάντηση. Π.χ. οι γωνίες του πίνακα, του τοίχου, του θρανίου.
\sol{13} Οι δύο ημιευθείες χωρίζουν το επίπεδο σε δύο μέρη. Κυρτή είναι αυτή που είναι το μικρότερο μέρος.
\sol{14} Από κάθε κορυφή φεύγουν $\nu-3$ διαγώνιοι (όχι προς τον εαυτό της και τις δύο γειτονικές). Σε όλες τις κορυφές $\nu(\nu-3)$, αλλά κάθε διαγώνιος μετράται δύο φορές, άρα $\frac{\nu(\nu-3)}{2}$. Στο δωδεκάγωνο $\frac{12\cdot9}{2}=54$.
\sol{15} Όχι: ένα ορθογώνιο $2\times6$ και ένα $3\times4$ έχουν ίδιο εμβαδόν αλλά δεν είναι ίσα.
\sol{16} Όχι με την έννοια του πολυγώνου που ξέρουμε, γιατί η γραμμή τέμνει τον εαυτό της και δεν ορίζει ένα εσωτερικό.
\sol{17} Ελεύθερη απάντηση.
\sol{18} Ελεύθερη απάντηση.
